% !TeX root = ../QFT_main.tex
\documentclass[../QFT_main.tex]{subfiles}
\begin{document}
%=========================================================
%=========================================================
\chapter{Finite Temperature QFT}\label{ch:finite_temperature_QFT}
%========================================================
%=========================================================


%---------------------------------------
%================================================
\section{Evolution in imaginary time}
%================================================
%---------------------------------------------



\begin{mytheorem}[$\tau$-evolution \& $\tau$-ordering]
Given a grand-canonical Hamiltonian $K=H-\mu N$, define the $\tau$-evolution of an operator and of a state by
{\small\begin{align}\begin{aligned}
    A_K(\tau)&:=e^{\tfrac{\tau}{\hbar}K}A e^{-\tfrac{\tau}{\hbar}K}\,,
    &
    |\chi\rangle_K(\tau)&:=e^{-\tfrac{\tau}{\hbar}K}|\chi\rangle\,.
\end{aligned}\end{align}}
Note the analogy with ordinary time evolution under the replacements
{\small\begin{align}\begin{aligned}
    H\longmapsto K=H-\mu N\,,
    \qquad
    t\longmapsto-i\tau\,,
    \qquad
    it\longmapsto\tau\,,
    \qquad
    i\hbar\partial_t=-\hbar\partial_{it}\longmapsto-\hbar\partial_\tau\,.
\end{aligned}\end{align}}
In other words, the time-evolution operator $e^{-\tfrac{i}{\hbar}tH}$ can be mapped to the statistical operator $e^{-\beta K}/Z$ (modulo the normalization $Z$) if we enrich the Hamiltonian with an external potential energy, namely the chemical-potential energy $-\mu N$, and identify time with imaginary inverse temperature.
Some remarks are in order.
\begin{itemize}
    \item If the Hamiltonian is $\tau$-dependent, we consider $U(\tau,0)$ instead of $e^{-\tfrac{\tau}{\hbar}K}$.
    \item We will typically consider evolution times $\tau\in[0,\beta\hbar]$ (or $\tau\in[-\beta\hbar/2,\beta\hbar/2]$ in the symmetric case) and reference time $\tau_0=0$.
\end{itemize}

The $\tau$-evolved operator satisfies the Heisenberg equation
{\small\begin{align}\begin{aligned}
    -\hbar\tfrac{\partial}{\partial\tau}A_K(\tau)=[A_K(\tau),K]\,.
\end{aligned}\end{align}}
If the basis $|a\rangle$ diagonalizes the Hamiltonian $H=\sum_a\epsilon_a c_a^\dagger c_a$ then we have
{\small\begin{align}\begin{aligned}
    c_{a,K}(\tau)&=e^{-\tfrac{\tau}{\hbar}(\epsilon_a-\mu)}c_a\,,
    &
    c_{a,K}^\dagger(\tau)&=e^{\tfrac{\tau}{\hbar}(\epsilon_a-\mu)}c_a^\dagger\,.
\end{aligned}\end{align}}
The $\tau$-evolution of a composition of operators is the composition of the $\tau$-evolved operators,
{\small\begin{align}\begin{aligned}
    (AB)_K(\tau)=A_K(\tau)B_K(\tau)\,.
\end{aligned}\end{align}}
%
The $\tau$-ordering of $\tau$-evolved operators is defined by linearity, starting from a basis of field operators $\psi_a$, e.g.\ the operators $c_a,c_b^\dagger$:
{\small\begin{align}\begin{aligned}
    T\big\{\psi_{1,K}(\tau_1)\cdots(\psi_{n,K})(\tau_n)\big\}
    =
    (\pm1)^p
    \psi_{p_1,K}(\tau_{p_1})\cdots\psi_{p_n,K}(\tau_{p_n})\,,
    \qquad
    \tau_{p_1}\geq\cdots\geq\tau_{p_n}\,,
\end{aligned}\end{align}}
where $p$ is the permutation that orders the times, with sign $\pm1$ for bosonic and fermionic operators, respectively.
This definition is then extended by linearity to general operators $O$ that are linear combinations of compositions of field operators.
\end{mytheorem}


\begin{mytheorem}[Evolution pictures]
The thermal Hamiltonian $K=H-\mu N$ is split as $K=K_0+V$, with $K_0$ a thermal Hamiltonian that we can treat and $V$ a (hopefully small) perturbation.
We fix a reference time $\tau_0$ at which the three evolution pictures coincide; in the thermal case we shall always take $\tau_0=0$.
The evolution operator in the Schr\"odinger picture satisfies
{\small\begin{align}\begin{aligned}
    \left\{
    \begin{aligned}
        -\hbar\tfrac{\partial}{\partial\tau}U(\tau,\sigma)&=K(\tau)U(\tau,\sigma)\,,
        \\
        U(\sigma,\sigma)&=\mathrm{id}\,,
    \end{aligned}
    \right.
\end{aligned}\end{align}}
with solution
{\small\begin{align}\begin{aligned}
    U(\tau,s)
    =
    T\Big\{
        \exp\big(-\tfrac{1}{\hbar}\int_s^\tau K(\varphi)\,d\varphi\big)
    \Big\}\,.
\end{aligned}\end{align}}
If $K$ is time-independent, this becomes simply $U(\tau,s)=e^{-\tfrac{\tau-s}{\hbar}K}$.

The Schr\"odinger propagator $U(\tau,s)$ and interaction-picture propagator $U_I(\tau,s)$ are related by
{\small\begin{align}\begin{aligned}
    U(\tau,s)
    &=
    e^{-\tfrac{\tau-\tau_0}{\hbar}K_0}
    \underbrace{
        e^{\tfrac{\tau-\tau_0}{\hbar}K_0}
        U(\tau,s)
        e^{-\tfrac{s-\tau_0}{\hbar}K_0}
    }_{=:U_I(\tau,s)}
    e^{\tfrac{s-\tau_0}{\hbar}K_0}
    \notag\\
    &=
    e^{-\tfrac{\tau-\tau_0}{\hbar}K_0}
    U_I(\tau,s)
    e^{\tfrac{s-\tau_0}{\hbar}K_0}\,.
\end{aligned}\end{align}}
The evolution operator in the interaction picture satisfies the Dyson equation
{\small\begin{align}\begin{aligned}
    \left\{
    \begin{aligned}
        -\hbar\tfrac{\partial}{\partial\tau}U_I(\tau,\sigma)
        &=V_{K_0}(\tau)U_I(\tau,\sigma)\,,
        \\
        U_I(\sigma,\sigma)&=\mathrm{id}\,,
    \end{aligned}
    \right.
\end{aligned}\end{align}}
with solution
{\small\begin{align}\begin{aligned}
    U_I(\tau,\sigma)
    =
    T\Big\{
        \exp\big(-\tfrac{1}{\hbar}\int_\sigma^\tau V_{K_0}(\varphi)\,d\varphi\big)
    \Big\}\,.
\end{aligned}\end{align}}

The evolution operators in any picture satisfy 
{\small\begin{align}\begin{aligned}
    U_\bullet(\tau,\sigma)^{-1}&=U_\bullet(\sigma,\tau)\,,
    &
    U_\bullet(\tau,\sigma)U_\bullet(\sigma,\rho)&=U_\bullet(\tau,\rho)
    \qquad\forall\,\sigma,\tau,\rho\,,
\end{aligned}\end{align}}
but $U_\bullet(\tau,s)$ is no longer self-adjoint, so that $U_\bullet^\dagger\neq U_\bullet^{-1}$.

In any picture, the evolution of states and operators starts from time $\tau_0$.
Hnece, an intrinsically $\tau$-dependent operator $V(\tau)$ evolves in the different pictures as
{\small\begin{align}\begin{aligned}
    \text{Schr\"odinger:}\quad&
    V(\tau)\,,
    &
    \text{interaction:}\quad&
    V_{K_0}(\tau)
    =
    e^{\tfrac{\tau-\tau_0}{\hbar}K_0}
    V(\tau)
    e^{-\tfrac{\tau-\tau_0}{\hbar}K_0}\,,
    \notag\\
    \text{Heisenberg:}\quad&
    V_{K(\tau)}(\tau)
    =
    U(\tau,\tau_0)^{-1}V(\tau)U(\tau,\tau_0)\,.
\end{aligned}\end{align}}
\end{mytheorem}


\begin{mytheorem}[Averages with $\tau$-dependent Hamiltonians \& symmetrized setting]
Thermal averages for a generic $\tau$-dependent Hamiltonian $K(\tau)$ are defined by
{\small\begin{align}\label{eq:thermal_avg_time_dep_hamiltonian}
\begin{aligned}
    \langle\widehat{O}\rangle
    &:=
    \tfrac{\operatorname{Tr}\big\{U(\beta\hbar,0)\widehat{O}\big\}}
    {\operatorname{Tr}\big\{U(\beta\hbar,0)\big\}}\,,
    &
    \rho&:=\tfrac{1}{Z}U(\beta\hbar,0)\,,
    &
    Z&:=\operatorname{Tr}\big\{U(\beta\hbar,0)\big\}\,.
\end{aligned}\end{align}}
For $\tau$-independent Hamiltonians everything reduces to $U(\beta\hbar,0)=e^{-\beta K}$.

Alternatively, we can consider a symmetrized setting, more suitable for the zero-temperature limit, so that
{\small\begin{align}\begin{aligned}
    U\big(\tfrac{\beta\hbar}{2},-\tfrac{\beta\hbar}{2}\big)
    \longrightarrow U(+\infty,-\infty)\qquad \text{as}\quad \beta \to +\infty(1+\epsilon)\,.
\end{aligned}\end{align}}
The reference coincidence time for the three pictures is still taken to be $\tau_0=0$.
Operators (and/or states) are now evolved between times in $[-\beta\hbar/2,\beta\hbar/2]$, always with the suitable propagator $U(\tau,s)$.
Thermal averages are now defined symmetrically by
{\small\begin{align}\label{eq:symmetric_thermal_avg}
\begin{aligned}
    \langle\widehat{O}\rangle
    &:=
    \tfrac{
        \operatorname{Tr}\left\{
            U\big(\tfrac{\beta\hbar}{2},0\big)
            \widehat{O}
            U\big(0,-\tfrac{\beta\hbar}{2}\big)
        \right\}
    }{
        \operatorname{Tr}\left\{
            U\big(\tfrac{\beta\hbar}{2},-\tfrac{\beta\hbar}{2}\big)
        \right\}
    }\,.
\end{aligned}\end{align}}
Correlators still depend on the time differences $s,\tau\in[-\beta\hbar/2,\beta\hbar/2]$, so that $\tau-s\in[-\beta\hbar,\beta\hbar]$.
All the results are unchanged, with the obvious formal modifications.
\end{mytheorem}


\begin{mytheorem}[Thermal averages vs.\ $\tau$-evolution]
%
Therma averages of $\tau$-evolved operators obey several properties.
\begin{itemize}
    \item Thermal averages of single operators, or of products of equal-time $\tau$-evolved operators, are independent of $\tau$-evolution if the Hamiltonian $K$ is $\tau$-independent:
    {\small\begin{align}\begin{aligned}
        \langle A_K(\tau)\rangle_K
        &=
        \tfrac{1}{Z}\operatorname{Tr}\left\{
            e^{-\beta K}
            e^{\tfrac{\tau}{\hbar}K}
            A
            e^{-\tfrac{\tau}{\hbar}K}
        \right\}
        =\langle A\rangle_K\,.
    \end{aligned}\end{align}}
    \end{itemize}
    %
    \item If $K$ is time-independent, correlators depend only on time differences:
    {\small\begin{align}\begin{aligned}
        \langle A_K(\tau)B_K(\tau')\rangle_K
        &=
        \langle A_K(\tau-\tau')B\rangle_K\,.
    \end{aligned}\end{align}}
    Since {\small$\tau,\tau'\in[0,\hbar\beta]$}, or {\small$\tau,\tau'\in[-\tfrac{\hbar\beta}{2},\tfrac{\hbar\beta}{2}]$} in the symmetrized setting, we have {\small$\tau-\tau'\in[-\hbar\beta,\hbar\beta]$}.
    Indeed, cyclicity of the trace gives
    {\small\begin{align}\begin{aligned}
        &\operatorname{Tr}\left\{
            \tfrac{e^{-\beta K}}{Z}
            e^{\tfrac{\tau}{\hbar}K}A e^{-\tfrac{\tau}{\hbar}K}
            e^{\tfrac{\tau'}{\hbar}K}B e^{-\tfrac{\tau'}{\hbar}K}
        \right\}
        =
        \operatorname{Tr}\left\{
            \tfrac{e^{-\beta K}}{Z}
            e^{\tfrac{\tau-\tau'}{\hbar}K}
            A e^{-\tfrac{\tau-\tau'}{\hbar}K}B
        \right\}\,.
    \end{aligned}\end{align}}
    %
    \item \textbf{KMS property.} For any operators $A_1,\ldots,A_n$ and any thermal Hamiltonian $K$, the Kubo--Martin--Schwinger property reads
{\small\begin{align}\begin{aligned}
    \left\langle
        A_{1,K}(\tau_1)\cdots A_{n,K}(\tau_n)
    \right\rangle_K
    &=
    \left\langle
        A_{n,K}(\tau_n+\beta\hbar)
        A_{1,K}(\tau_1)\cdots A_{n-1,K}(\tau_{n-1})
    \right\rangle_K\,,
    \quad \tau_1,\ldots,\tau_n\in[0,\beta\hbar]\,.
\end{aligned}\end{align}}
This follows directly from cyclicity of the trace. Indeed, considerint two operators for simplicity, we have 
{\small\begin{align}\begin{aligned}
    \langle A_K(\tau)B_K(\sigma)\rangle_K
    &=
    \operatorname{Tr}\left\{
        \tfrac{e^{-\beta K}}{Z}A_K(\tau)B_K(\sigma)
    \right\}
    =
    \operatorname{Tr}\left\{
        B_K(\sigma)\tfrac{e^{-\beta K}}{Z}A_K(\tau)
    \right\}
    =
    \operatorname{Tr}\bigg\{
        \tfrac{e^{-\beta K}}{Z}
        \underbrace{
            e^{\beta K}B_K(\sigma)e^{-\beta K}
        }_{=B_K(\sigma+\beta\hbar)}
        A_K(\tau)
    \bigg\}\,.
\end{aligned}\end{align}}
\end{mytheorem}


\begin{mytheorem}[KMS property iff thermal Gibbs state]
%
The KMS property actually fully identifies thermal Gibbs states, i.e. an operator $\rho$ satisfies the KMS property if and only if it is a Gibbs thermal state.
Indeed, suppose that a state $\rho$ satisfies
{\small\begin{align}\begin{aligned}
    \operatorname{Tr}\{\rho A(\tau)B(\tau')\}
    &=
    \operatorname{Tr}\{\rho B(\tau'+\beta\hbar)A(\tau)\}, \quad \forall A,B,\tau,\tau'\,.
\end{aligned}\end{align}}
for every pair of operators $A,B$ and every pair of parameters $\tau,\tau'$.
Setting $\tau=\tau'=0$ gives
{\small\begin{align}\begin{aligned}
    \operatorname{Tr}\{\rho AB\}
    &=
    \operatorname{Tr}\{\rho B(\beta\hbar)A\}\,,
    \implies
    \operatorname{Tr}\left\{
        \left(B\rho-\rho B(\beta\hbar)\right)A
    \right\}
    =0, \quad \forall A,B\,.
\end{aligned}\end{align}}
Taking $A=B\rho-\rho B(\beta\hbar)$, the non-degeneracy of the Hilbert--Schmidt inner product $\langle \mathcal{O}_1,\mathcal{O}_2\rangle_{\text{HS}}:=\operatorname{Tr}\{ \mathcal{O}_1^\dagger \mathcal{O}_2 \}$ implies 
{\small\begin{align}\begin{aligned}
    0=B\rho-\rho B(\beta\hbar)=\big(B\rho e^{\beta K}-\rho e^{\beta K}B\big)e^{-\beta K}
    \quad \text{that is}\quad 
    [B,\rho e^{\beta K}]&=0, \quad\forall B\,. 
\end{aligned}\end{align}}
Since this holds for every operator $B$, we conclude that $\rho e^{\beta K}$ is a multiple of the identity, i.e. $\rho = \lambda e^{-\beta K}$ for some $\lambda\in\mathbb{C}$.
Enforcing self-adjointness and normalization gives $\lambda=1/Z$ with $Z=\operatorname{Tr}\{e^{-\beta K}\}$.
\end{mytheorem}



\begin{mytheorem}[Thermal averages of products of creation/destruction operators \& symmetries of the Hamiltonian (Molinari 2-2020 -- 4AB)]
Let $S$ be a one-body operator which is a symmetry, or conserved quantity, of the thermal Hamiltonian:
{\small\begin{align}\begin{aligned}
    [K,S]=0\,.
\end{aligned}\end{align}}
We can find a single-particle basis $|s\rangle$ diagonalizing $S$, so that in second quantization
{\small\begin{align}\begin{aligned}
    S=\sum_s\langle s|S|s\rangle c_s^\dagger c_s
    =\sum_s s\,c_s^\dagger c_s\,.
\end{aligned}\end{align}}
In this basis,
{\small\begin{align}\begin{aligned}
    e^{\lambda S}c_s e^{-\lambda S}&=e^{-\lambda s}c_s\,,
    &
    e^{\lambda S}c_s^\dagger e^{-\lambda S}&=e^{\lambda s}c_s^\dagger\,.
\end{aligned}\end{align}}
Consider a product $\widetilde c_1\cdots\widetilde c_k$ of creation and destruction operators.
If $\mathcal C$ and $\mathcal A$ denote the sets of creators and annihilators appearing in the product, respectively, then
{\small\begin{align}\begin{aligned}
    e^{\lambda S}
    \widetilde c_1\cdots\widetilde c_k
    e^{-\lambda S}
    &=
    e^{\lambda\left(\sum_{j\in\mathcal C}s_j-\sum_{j\in\mathcal A}s_j\right)}
    \widetilde c_1\cdots\widetilde c_k\,,
    \\
    [S,\widetilde c_1\cdots\widetilde c_k]
    &=
    \left(\sum_{j\in\mathcal C}s_j-\sum_{j\in\mathcal A}s_j\right)
    \widetilde c_1\cdots\widetilde c_k\,.
\end{aligned}\end{align}}
Taking the thermal average and using $[S,K]=0$ together with cyclicity of the trace gives
{\small\begin{align}\begin{aligned}
    \left(\sum_{j\in\mathcal C}s_j-\sum_{j\in\mathcal A}s_j\right)
    \left\langle\widetilde c_1\cdots\widetilde c_k\right\rangle_K
    &=
    \left\langle[S,\widetilde c_1\cdots\widetilde c_k]\right\rangle_K
    =0\,.
\end{aligned}\end{align}}
Therefore,
{\small\begin{align}\begin{aligned}
    \left\langle\widetilde c_1\cdots\widetilde c_k\right\rangle_K=0
    \qquad\text{if}\qquad
    \sum_{j\in\mathcal C}s_j-\sum_{j\in\mathcal A}s_j\neq0\,.
\end{aligned}\end{align}}

Two important examples are the following.
\begin{itemize}
    \item \textbf{Conservation of the total particle number.}
    If $[K,N]=0$, where
    {\small\begin{align}\begin{aligned}
        N=\sum_r c_r^\dagger c_r\,,
    \end{aligned}\end{align}}
    then the argument works in any single-particle basis and for any product containing both creation and destruction operators.
    Indeed,
    {\small\begin{align}\begin{aligned}
        e^{-\lambda N}c_r e^{\lambda N}&=e^\lambda c_r\,,
        &
        e^{-\lambda N}c_r^\dagger e^{\lambda N}&=e^{-\lambda}c_r^\dagger\,,
    \end{aligned}\end{align}}
    and hence the thermal average of a product of field operators vanishes unless the number of creators equals the number of annihilators:
    {\small\begin{align}\begin{aligned}
        0
        &=
        \operatorname{Tr}\left\{\tfrac{e^{-\beta K}}{Z}
        [N,\widetilde c_1\cdots\widetilde c_k]\right\}
        \\
        &=
        \left(\#\text{ creators}-\#\text{ annihilators}\right)
        \left\langle\widetilde c_1\cdots\widetilde c_k\right\rangle_K\,.
    \end{aligned}\end{align}}
    In particular, such averages always vanish for an odd number of field operators.
    For example, $\langle c_k^\dagger\psi(x)c_p\rangle_K=0$ after expanding $\psi(x)$ in terms of annihilation operators.

    \item \textbf{Conservation of the total momentum.}
    If $[K,P]=0$, we use the momentum basis in which
    {\small\begin{align}\begin{aligned}
        P=\sum_r p_r c_r^\dagger c_r\,,
    \end{aligned}\end{align}}
    where $p_r$ is the momentum carried by the state $|r\rangle$.
    Then
    {\small\begin{align}\begin{aligned}
        [P,\widetilde c_1\cdots\widetilde c_k]
        &=
        \left(\sum_{j\in\mathcal C}p_j-\sum_{j\in\mathcal A}p_j\right)
        \widetilde c_1\cdots\widetilde c_k\,,
        \\
        0
        &=
        \left(\sum_{j\in\mathcal C}p_j-\sum_{j\in\mathcal A}p_j\right)
        \left\langle\widetilde c_1\cdots\widetilde c_k\right\rangle_K\,.
    \end{aligned}\end{align}}
    Thus, thermal averages of products of field operators vanish unless the total momentum of the creators equals the total momentum of the annihilators.
\end{itemize}

\emph{Why this works.}
The argument above avoids the difficulty of finding an explicit common basis of eigenvectors of $S$ and $K$ in the Fock space built from the single-particle states $|s\rangle$.
Such a common basis always exists because $[K,S]=0$, although it need not be manifestly of the form $|s_1,\ldots,s_n\rangle$.
For the sake of the heuristic, suppose that we have found such a basis.
A state
{\small\begin{align}\begin{aligned}
    |s_1,\ldots,s_n\rangle
    \propto c_{s_1}^\dagger\cdots c_{s_n}^\dagger|0\rangle
\end{aligned}\end{align}}
is an eigenvector of $S$ with eigenvalue $s_1+\cdots+s_n$.
Inserting this basis in the trace gives
{\small\begin{align}\begin{aligned}
    \left\langle\widetilde c_1\cdots\widetilde c_k\right\rangle_K
    &=
    \sum_\alpha
    \langle\alpha|\tfrac{e^{-\beta K}}{Z}|\alpha\rangle
    \langle\alpha|\widetilde c_1\cdots\widetilde c_k|\alpha\rangle\,.
\end{aligned}\end{align}}
When $\sum_{j\in\mathcal C}s_j-\sum_{j\in\mathcal A}s_j\neq0$, the product maps $|\alpha\rangle$ to an orthogonal eigenspace of $S$, and every diagonal matrix element in this sum vanishes.
\end{mytheorem}







\begin{mytheorem}[Matsubara frequencies for $\tau$-ordered correlators (Molinari 2-2020 -- 5A)]
Define the $\tau$-ordered correlator by
{\small\begin{align}\begin{aligned}
    -C_{AB}^{T}(\tau-\sigma)
    &:=
    \left\langle
        T\{A_K(\tau)B_K(\sigma)\}
    \right\rangle_K
    \\
    &=
    -\sum_{n\in\mathbb Z}
    \tfrac{1}{\beta\hbar}
    e^{-i\omega_n(\tau-\sigma)}
    \widetilde C_{AB}^{T}(i\omega_n)\,,
    \qquad
    \omega_n=\tfrac{n\pi}{\beta\hbar}\,.
\end{aligned}\end{align}}
Its Matsubara coefficients are
{\small\begin{align}\begin{aligned}
    \widetilde C_{AB}^{T}(i\omega_n)
    &=
    \begin{cases}
        \displaystyle
        \int_0^{\beta\hbar}
        e^{i\omega_n\rho}C_{AB}^{T}(\rho)\,d\rho\,,
        &
        \begin{array}{l}
            n\text{ even if }A,B\text{ commute under }\tau\text{-ordering},\\
            n\text{ odd if }A,B\text{ anticommute under }\tau\text{-ordering},
        \end{array}
        \\[1.1em]
        0\,,
        &
        \begin{array}{l}
            n\text{ odd in the commuting case},\\
            n\text{ even in the anticommuting case}.
        \end{array}
    \end{cases}
\end{aligned}\end{align}}
Thus density--density correlators, such as $\langle T\{n(x)n(y)\}\rangle_K$, carry even or bosonic Matsubara frequencies, while fermionic Green functions, such as $\langle T\{\psi(x)\psi^\dagger(y)\}\rangle_K$, carry odd or fermionic Matsubara frequencies.

The relativistic symmetric convention for the Fourier transform on $[-\beta\hbar,\beta\hbar]$ is
{\small\begin{align}\begin{aligned}
    f(t)
    &=
    \sum_{n\in\mathbb Z}
    \tfrac{1}{\beta\hbar}
    f(i\omega_n)e^{-i\omega_n t}\,,
    &
    f(i\omega_n)
    &=
    \tfrac12\int_{-\beta\hbar}^{\beta\hbar}
    f(t)e^{i\omega_n t}\,dt\,.
\end{aligned}\end{align}}
Operators can be evolved at times $\tau,s\in[0,\beta\hbar]$, or at times $\tau,s\in[-\beta\hbar/2,\beta\hbar/2]$ in the symmetric setting.
In either convention, their correlators depend on time differences in $[-\beta\hbar,\beta\hbar]$.
Expanding on the corresponding Fourier basis extends them $2\beta\hbar$-periodically outside this interval.

\emph{Proof.}
For $0\leq\tau,\tau'\leq\beta\hbar$, the correlator is a function of $\tau-\tau'\in[-\beta\hbar,\beta\hbar]$.
Its Fourier coefficient on the symmetric interval is
{\small\begin{align}\begin{aligned}
    C_{AB}(i\omega_n)
    &=
    \tfrac12\int_{-\beta\hbar}^{\beta\hbar}
    e^{i\omega_n\sigma}C_{AB}^{T}(\sigma)\,d\sigma
    \\
    &=
    \tfrac12\int_0^{\beta\hbar}
    e^{i\omega_n\sigma}
    \left[
        C_{AB}^{T}(\sigma)
        +(-1)^n C_{AB}^{T}(\sigma-\beta\hbar)
    \right]d\sigma\,.
\end{aligned}\end{align}}
For $0\leq\sigma\leq\beta\hbar$, the KMS rule and $\tau$-ordering imply
{\small\begin{align}\begin{aligned}
    -C_{AB}^{T}(\sigma-\beta\hbar)
    &=
    \left\langle T\{A_K(\sigma-\beta\hbar)B\}\right\rangle_K
    \\
    &=
    \pm\left\langle BA_K(\sigma-\beta\hbar)\right\rangle_K
    =
    \pm\left\langle A_K(\sigma)B\right\rangle_K
    =
    \mp C_{AB}^{T}(\sigma)\,,
\end{aligned}\end{align}}
where the upper sign applies to commuting operators and the lower sign to anticommuting operators.
Therefore,
{\small\begin{align}\begin{aligned}
    C_{AB}(i\omega_n)
    &=
    \tfrac12\left[1\pm(-1)^n\right]
    \int_0^{\beta\hbar}
    e^{i\omega_n\sigma}C_{AB}^{T}(\sigma)\,d\sigma\,.
\end{aligned}\end{align}}
For commuting operators only even $n$ survive; for anticommuting operators only odd $n$ survive.
\end{mytheorem}


\begin{mytheorem}[Time-homogeneous systems vs.\ Matsubara frequencies]
Matsubara frequencies are a standard Fourier transform for imaginary times in $[-\beta\hbar,\beta\hbar]$.
For time-homogeneous systems, i.e.\ systems with a time-independent Hamiltonian, every correlator depends only on time differences.
Passing to Matsubara frequencies therefore turns convolution integrals into ordinary products:
{\small\begin{align}\begin{aligned}
    \int d\tau\,f(s-\tau)g(\tau-\rho)
    \quad\longleftrightarrow\quad
    f(i\omega_n)g(i\omega_n)\,.
\end{aligned}\end{align}}

For example, the time-dependent part of a proper self-energy insertion contains
{\small\begin{align}\begin{aligned}
    I(t_1-t_2)
    &:=
    \int_0^{\beta\hbar}dt_3
    \int_0^{\beta\hbar}dt_4\,
    G(t_1-t_3)
    \Sigma^\star(t_3-t_4)
    G^0(t_4-t_2)
    \\
    &=
    \sum_{n,m,q}
    \tfrac{
        \widehat G(i\omega_n)
        \widehat\Sigma^\star(i\omega_m)
        \widehat G^0(i\omega_q)
    }{(\beta\hbar)^3}
    \\
    &\qquad\qquad\times
    \int_0^{\beta\hbar}dt_3
    \int_0^{\beta\hbar}dt_4\,
    e^{i\omega_n(t_1-t_3)}
    e^{i\omega_m(t_3-t_4)}
    e^{i\omega_q(t_4-t_2)}
    \\
    &=
    \sum_n\tfrac{1}{\beta\hbar}
    e^{i\omega_n(t_1-t_2)}
    \widehat G(i\omega_n)
    \widehat\Sigma^\star(i\omega_n)
    \widehat G^0(i\omega_n)\,.
\end{aligned}\end{align}}
The two time integrals collapse the three frequencies by orthogonality:
{\small\begin{align}\begin{aligned}
    \int_0^{\beta\hbar}
    \tfrac{dt}{\beta\hbar}
    e^{it(\omega_m-\omega_n)}
    &=
    \begin{cases}
        1\,,&\omega_m=\omega_n\,,\\[0.4em]
        \displaystyle
        \tfrac{
            e^{i(\omega_m-\omega_n)\beta\hbar}-1
        }{
            i(\omega_m-\omega_n)\beta\hbar
        }
        =0\,,&\omega_m\neq\omega_n\,.
    \end{cases}
\end{aligned}\end{align}}
The last equality holds for both bosonic and fermionic Matsubara frequencies because the difference of two frequencies of the same statistics is always an even multiple of $\pi/(\beta\hbar)$.
\end{mytheorem}


\begin{mytheorem}[Matsubara coefficients for Dirac deltas of static interactions]
For a static two-body potential, the perturbative expansion of the evolution operator contains
{\small\begin{align}\begin{aligned}
    \mathcal U(\beta\hbar,0)
    &=
    T\exp\left[
        -\tfrac{1}{2\hbar}
        \int dx\,dx'\,
        U^0(x,x')
        \psi^\dagger(x)
        \psi^\dagger(x')
        \psi(x')
        \psi(x)
    \right]\,,
\end{aligned}\end{align}}
where
{\small\begin{align}\begin{aligned}
    U(x,x')=v(x,x')\delta(\tau-\tau')\,.
\end{aligned}\end{align}}
The interaction $V_{K_0}(\tau)$ is a function on $[0,\beta\hbar]$, and the delta distribution sets equal the two times of the instantaneous interaction.
It must therefore be expanded on the Fourier basis of $[0,\beta\hbar]$.
The orthonormal basis is
{\small\begin{align}\begin{aligned}
    e_n(\tau)
    &=
    \tfrac{1}{\sqrt{\beta\hbar}}e^{i\nu_n\tau}\,,
    &
    \nu_n&=\tfrac{2\pi n}{\beta\hbar}\,,
    \qquad n\in\mathbb Z\,.
\end{aligned}\end{align}}
These are precisely the bosonic frequencies, i.e.\ the even members of $\omega_n=n\pi/(\beta\hbar)$ on the symmetric interval.
Since
{\small\begin{align}\begin{aligned}
    \int_0^{\beta\hbar}e_n(\tau)\delta(\tau)\,d\tau
    &=
    \tfrac{1}{\sqrt{\beta\hbar}}\,,
\end{aligned}\end{align}}
we obtain
{\small\begin{align}\begin{aligned}
    \delta(\tau-\tau')
    &=
    \sum_{n\in\mathbb Z}
    \tfrac{1}{\sqrt{\beta\hbar}}e_n(\tau-\tau')
    \\
    &=
    \tfrac{1}{\beta\hbar}
    \sum_{\substack{n\in\mathbb Z\\ n\ \mathrm{even}}}
    e^{i\omega_n(\tau-\tau')}\,,
    \qquad
    -\beta\hbar\leq\tau-\tau'\leq\beta\hbar\,.
\end{aligned}\end{align}}
Thus the delta distribution on $[0,\beta\hbar]$ has the same expansion as a bosonic correlator.
This is naturally related to the fact that a static interaction can arise as the static limit of a contraction, or correlator, of two bosons mediating the interaction.

There is a distinction between the delta distribution on $[0,\beta\hbar]$ and the one on the doubled interval $[-\beta\hbar,\beta\hbar]$:
{\small\begin{align}\begin{aligned}
    \delta_{[0,\beta\hbar]}(\tau-\tau')
    &=
    \tfrac{1}{\beta\hbar}
    \sum_{\substack{n\in\mathbb Z\\ n\ \mathrm{even}}}
    e^{i\omega_n(\tau-\tau')}\,,
    \\
    \delta_{[-\beta\hbar,\beta\hbar]}(\tau-\tau')
    &=
    \tfrac{1}{2\beta\hbar}
    \sum_{n\in\mathbb Z}
    e^{i\omega_n(\tau-\tau')}\,.
\end{aligned}\end{align}}
Testing the two distributions against functions on their respective domains reproduces the two Fourier series:
{\small\begin{align}\begin{aligned}
    f(\tau)
    &=
    \sum_{n\in\mathbb Z}
    \tfrac{1}{\beta\hbar}
    \left[
        \int_0^{\beta\hbar}
        e^{i\nu_n\varphi}f(\varphi)\,d\varphi
    \right]
    e^{-i\nu_n\tau}\,,
    \\
    g(\tau)
    &=
    \sum_{n\in\mathbb Z}
    \tfrac{1}{2\beta\hbar}
    \left[
        \int_{-\beta\hbar}^{\beta\hbar}
        e^{i\omega_n\varphi}g(\varphi)\,d\varphi
    \right]
    e^{-i\omega_n\tau}\,.
\end{aligned}\end{align}}
\end{mytheorem}


\begin{mytheorem}[General method for Matsubara sums (Molinari 2-2020 -- 5B, 6B)]
Introduce the functions related to the Bose and Fermi distributions,
{\small\begin{align}\begin{aligned}
    n_{\mp}(z)
    &=
    \tfrac{1}{e^{\beta\hbar z}\mp1}\,.
\end{aligned}\end{align}}
They have simple poles at $z_n=i\omega_n$, with $n$ even for $n_-$ and $n$ odd for $n_+$.
Their residues are, respectively, $\pm1/(\beta\hbar)$.

Let $f$ be a meromorphic function which decays at infinity and whose poles $\{z_p\}$ differ from those of $n_-$ or $n_+$.
For a small convergence parameter $\eta$, consider the integral on a large circle $C$:
{\small\begin{align}\begin{aligned}
    \oint_C\tfrac{dz}{2\pi i}\,
    e^{\eta z}n_{\mp}(z)f(z)
    &=
    \pm\tfrac{1}{\beta\hbar}
    \sum_n f(i\omega_n)e^{i\omega_n\eta}
    +
    \sum_p\operatorname{Res}\left(f n_{\mp},z_p\right)\,.
\end{aligned}\end{align}}
The factor $e^{\eta z}$, with $\eta$ ultimately taken to zero, ensures convergence on the half-circle with $\operatorname{Re}z<0$, while $n_-$ decays for $\operatorname{Re}z>0$.
Since the contour integral vanishes at infinite radius, the Matsubara sum is
{\small\begin{align}\begin{aligned}
    \tfrac{1}{\beta\hbar}
    \sum_n f(i\omega_n)e^{i\omega_n\eta}
    &=
    \mp\sum_p\operatorname{Res}\left(f n_{\mp},z_p\right)\,.
\end{aligned}\end{align}}
\end{mytheorem}


\begin{mytheorem}[Matsubara sum for the density]
Taking
{\small\begin{align}\begin{aligned}
    f(z)
    &=
    \tfrac{1}{
        z-\tfrac{1}{\hbar}(\epsilon_r-\mu)
    }
\end{aligned}\end{align}}
in the residue formula gives
{\small\begin{align}\begin{aligned}
    \tfrac{1}{\beta\hbar}
    \sum_{\omega_n}
    \tfrac{e^{i\omega_n\eta}}{
        i\omega_n-\tfrac{1}{\hbar}(\epsilon_r-\mu)
    }
    &=
    \mp\tfrac{1}{
        e^{\beta(\epsilon_r-\mu)}\mp1
    }
    =\mp n_r\,.
\end{aligned}\end{align}}

There is also an alternative proof which avoids complex contour integrals.
For a field component $\psi_\mu(x)$,
{\small\begin{align}\begin{aligned}
    \langle n_\mu(x)\rangle
    &:=
    \langle\psi_\mu^\dagger(x)\psi_\mu(x)\rangle
    \\
    &=
    \pm\left\langle
        T\{\psi_\mu(x,\tau)\psi_\mu^\dagger(x,\tau^+)\}
    \right\rangle
    =
    \mp G(x,\tau;x,\tau^+)
    \\
    &=
    \mp\tfrac{1}{\beta\hbar}
    \sum_{\omega_n}
    e^{-i\omega_n(\tau-\tau^+)}
    G(x\mu,x\mu;i\omega_n)
    \\
    &=
    \mp\sum_r|\langle x\mu|r\rangle|^2
    \tfrac{1}{\beta\hbar}
    \sum_{\omega_n}
    \tfrac{e^{i\omega_n\eta}}{
        i\omega_n-\tfrac{1}{\hbar}(\epsilon_r-\mu)
    }\,.
\end{aligned}\end{align}}
On the other hand, expanding the fields in the basis $|r\rangle$ gives
{\small\begin{align}\begin{aligned}
    \langle n_\mu(x)\rangle
    &=
    \sum_{r,r'}
    \langle r|x\mu\rangle
    \langle x\mu|r'\rangle
    \langle c_r^\dagger c_{r'}\rangle_K
    \\
    &=
    \sum_{r,r'}
    \langle r|x\mu\rangle
    \langle x\mu|r'\rangle
    \delta_{rr'}n_r
    =
    \sum_r|\langle x\mu|r\rangle|^2n_r\,.
\end{aligned}\end{align}}
Comparing the two expressions gives the result.
More directly, in the basis diagonalizing $K$,
{\small\begin{align}\begin{aligned}
    n_r
    &=
    \langle\widehat n_r\rangle
    =
    \langle c_r^\dagger c_r\rangle
    =
    \mp G(r,\tau;r,\tau^+)
    \\
    &=
    \mp\tfrac{1}{\beta\hbar}
    \sum_n
    \tfrac{e^{i\omega_n\eta}}{
        i\omega_n-\tfrac{1}{\hbar}(\epsilon_r-\mu)
    }\,.
\end{aligned}\end{align}}
\end{mytheorem}


\begin{mytheorem}[Other Matsubara sums \& alternative methods]
Useful thermal series to evaluate with the residue method include
{\small\begin{align}\begin{aligned}
    &\tfrac{1}{\beta\hbar}
    \sum_n
    \tfrac{1}{
        i\omega_n-\tfrac{1}{\hbar}(\epsilon-\mu)
    }
    \tfrac{1}{
        i(\omega_n-\nu)-\tfrac{1}{\hbar}(\epsilon'-\mu)
    }\,,
    \\
    &\tfrac{1}{\beta\hbar}
    \sum_n
    \tfrac{1}{
        \left(
            i\omega_n-\tfrac{1}{\hbar}(\epsilon-\mu)
        \right)^2
    }\,,
    \\
    &\tfrac{1}{\beta\hbar}
    \sum_n e^{i\omega_n\eta}
    \log\left(
        i\omega_n-\tfrac{\epsilon-\mu}{\hbar}
    \right)\,.
\end{aligned}\end{align}}
The logarithmic sum requires a keyhole contour.

Some series can be summed more simply by partial fractions.
For example,
{\small\begin{align}\begin{aligned}
    \sum_n
    \tfrac{1}{i\omega_n-x}
    \tfrac{1}{i\omega_n-i\nu-y}
    &=
    \lim_{\eta\to0}
    \sum_n
    \left(
        \tfrac{e^{i\omega_n\eta}}{i\omega_n-x}
        -
        \tfrac{e^{i\omega_n\eta}}{i\omega_n-i\nu-y}
    \right)
    \tfrac{1}{x-i\nu-y}\,.
\end{aligned}\end{align}}
The two sums on the right are of the density type.
The convergence factor $e^{i\omega_n\eta}$ must be retained while evaluating them separately because each behaves as $1/\omega_n$, even though the original sum behaves as $1/\omega_n^2$ and converges.
This does not alter the result, since the left-hand side is convergent and $\lim_{\eta\to0}e^{i\omega_n\eta}=1$.
\end{mytheorem}







%---------------------------------------
%================================================
\section{Perturbative approach}
%================================================
%---------------------------------------------

\begin{mytheorem}[Reduction formula, i.e.\ finite-temperature Gell-Mann--Low formula]
%
Split the thermal Hamiltonian as $K=K_0+V$.
For any operators $A_1,\ldots,A_n$, the reduction formula reads
{\small\begin{align}\begin{aligned}
    &\Big\langle
        T\big\{
            A_{1,K}(\tau_1)\cdots A_{n,K}(\tau_n)
        \big\}
    \Big\rangle_K
=
    \tfrac{
        \Big\langle
            T\big\{
                U_I(\beta\hbar,0)
                A_{1,K_0}(\tau_1)\cdots A_{n,K_0}(\tau_n)
            \big\}
        \Big\rangle_{K_0}
    }{
        \big\langle U_I(\beta\hbar,0)\big\rangle_{K_0}
    }\,.
\end{aligned}\end{align}}
Unlike the zero-temperature formula, this result requires no further assumptions concerning energy-level crossings.
\begin{itemize}
    \item For a $\tau$-dependent Hamiltonian the reduction formula holds identically, with thermal averages defined by \eqref{eq:thermal_avg_time_dep_hamiltonian}.
\item In the $\tau$-symmetric convention, including for $\tau$-dependent Hamiltonians, the obvious modification is
{\small\begin{align}\begin{aligned}
    &\Big\langle
        T\big\{
            A_{1,K}(\tau_1)\cdots A_{n,K}(\tau_n)
        \big\}
    \Big\rangle_K
=
    \tfrac{
        \Big\langle
            T\big\{
                U_I\big(\tfrac{\beta\hbar}{2},-\tfrac{\beta\hbar}{2}\big)
                A_{1,K_0}(\tau_1)\cdots A_{n,K_0}(\tau_n)
            \big\}
        \Big\rangle_{K_0}
    }{
        \Big\langle
            U_I\big(\tfrac{\beta\hbar}{2},-\tfrac{\beta\hbar}{2}\big)
        \Big\rangle_{K_0}
    }\,.
\end{aligned}\end{align}}
The proof is unchanged after recalling that symmetric thermal averages are defind by \eqref{eq:symmetric_thermal_avg}.
\end{itemize}
\end{mytheorem}

\begin{proof}
Assume first that $\tau_1\geq\cdots\geq\tau_n$.
Writing
{\small\begin{align}\begin{aligned}
    Q&:=\operatorname{Tr}\big\{e^{-\beta K}\big\}\,,
    &
    Q_0&:=\operatorname{Tr}\big\{e^{-\beta K_0}\big\}\,,
\end{aligned}\end{align}}
and repeatedly inserting the interaction-picture propagators gives
{\small\begin{align}\begin{aligned}
    \Big\langle
        T\big\{
            A_{1,K}(\tau_1)\cdots A_{n,K}(\tau_n)
        \big\}
    \Big\rangle_K
    &=
    \tfrac{Q_0}{Q}
    \Big\langle
        U_I(\beta\hbar,\tau_1)
        A_{1,K_0}(\tau_1)
        U_I(\tau_1,\tau_2)
        A_{2,K_0}(\tau_2)
        \cdots
        U_I(\tau_n,0)
    \Big\rangle_{K_0}
    \\
    &=
    \tfrac{Q_0}{Q}
    \Big\langle
        T\big\{
            U_I(\beta\hbar,0)
            A_{1,K_0}(\tau_1)\cdots A_{n,K_0}(\tau_n)
        \big\}
    \Big\rangle_{K_0}\,.
\end{aligned}\end{align}}
The factors in the middle line follow from
{\small\begin{align}\begin{aligned}
    U_I(\tau,s)
    &=
    e^{\tfrac{\tau}{\hbar}K_0}
    U(\tau,s)
    e^{-\tfrac{s}{\hbar}K_0}\,,
    &
    A_{i,K_0}(\tau_i)
    &=
    e^{\tfrac{\tau_i}{\hbar}K_0}
    A_i
    e^{-\tfrac{\tau_i}{\hbar}K_0}\,.
\end{aligned}\end{align}}
Taking $n=0$ yields
{\small\begin{align}\begin{aligned}
    1
    &=
    \langle 1 \rangle_K
    =\operatorname{Tr}\left\{e^{-\beta K}/Q\right\}_K
    =
    \tfrac{Q_0}{Q}
    \big\langle T\big\{ U_I(\beta\hbar,0)\big\}\big\rangle_{K_0}\,,
    \quad\text{so that}\quad
    \tfrac{Q_0}{Q}
    =
    \big\langle U_I(\beta\hbar,0)\big\rangle_{K_0}^{-1}\,,
\end{aligned}\end{align}}
which proves the formula.
If the times are not already ordered, both sides acquire the same factor $(\pm1)^p$ under the permutation that orders them, and these factors cancel.
\end{proof}


\begin{mytheorem}[Computing the full grand potential $\Omega$ in the interaction picture]
%
We can relate the free and interacting grand potentials in several ways, and then compute the relation perturbatively 
{\small\begin{align}\begin{aligned}
    \Omega&:=-\tfrac{1}{\beta}\log Q\,,
    &
    \Omega_0&:=-\tfrac{1}{\beta}\log Q_0\,\,.
\end{aligned}\end{align}}

\noindent
\textbf{1. Reduction formula.}
The $n=0$ case of the reduction formula gives
{\small\begin{align}\begin{aligned}
    \tfrac{Q}{Q_0}
    &=
    \big\langle U_I(\beta\hbar,0)\big\rangle_{K_0}\,,
    \implies 
    \Omega-\Omega_0= -\tfrac{1}{\beta}
    \log\big\langle U_I(\beta\hbar,0)\big\rangle_{K_0}\,.
\end{aligned}\end{align}}
The right-hand side can be computed perturbatively.
Expanding the interaction-picture propagator to second order gives
{\small\begin{align}\begin{aligned}
    \big\langle U_I(\beta\hbar,0)\big\rangle_{K_0}
    &=
    \Big\langle
        T\Big\{
            1
            -\tfrac{1}{\hbar}\int_0^{\beta\hbar}
                V_{K_0}(\sigma)\,d\sigma
            +\tfrac{1}{2\hbar^2}
                \int_0^{\beta\hbar}d\sigma_1
                \int_0^{\beta\hbar}d\sigma_2\,
                V_{K_0}(\sigma_1)V_{K_0}(\sigma_2)
            +\mathcal O(V^3)
        \Big\}
    \Big\rangle_{K_0}
    \\
    &=
    1-\beta\langle V\rangle_{K_0}
    +\tfrac{1}{2\hbar^2}
        \int_0^{\beta\hbar}d\sigma_1
        \int_0^{\beta\hbar}d\sigma_2\,
        \Big\langle
            T\big\{
                V_{K_0}(\sigma_1)V_{K_0}(\sigma_2)
            \big\}
        \Big\rangle_{K_0}
    +\mathcal O(V^3)\,,
\end{aligned}\end{align}}
where $\langle V_{K_0}(\sigma)\rangle_{K_0}=\langle V\rangle_{K_0}$ is independent of $\sigma$.
Using $\log(1+x)=x-\tfrac{x^2}{2}+\cdots$, we find
{\small\begin{align}\begin{aligned}
    \log\big\langle U_I(\beta\hbar,0)\big\rangle_{K_0}
    &=
    -\beta\langle V\rangle_{K_0}
    +\tfrac{1}{2\hbar^2}
        \int_0^{\beta\hbar}d\sigma_1
        \int_0^{\beta\hbar}d\sigma_2\,
        \Big\langle
            T\big\{
                V_{K_0}(\sigma_1)V_{K_0}(\sigma_2)
            \big\}
        \Big\rangle_{K_0}
    -\tfrac12\beta^2\langle V\rangle_{K_0}^2
    +\mathcal O(V^3)\,.
\end{aligned}\end{align}}
In conclusion, denoting $\delta V_{K_0}(\sigma):=V_{K_0}(\sigma)-\langle V\rangle_{K_0}$, the \textbf{grand potential offset is}
{\small\begin{align}\begin{aligned}
    \Omega-\Omega_0
    &=
    \langle V\rangle_{K_0}
    -\tfrac{1}{2\beta\hbar^2}
        \int_0^{\beta\hbar}d\sigma_1
        \int_0^{\beta\hbar}d\sigma_2\,
        \Big\langle
            T\big\{
                \delta V_{K_0}(\sigma_1)
                \delta V_{K_0}(\sigma_2)
            \big\}
        \Big\rangle_{K_0}
    +\mathcal O(V^3)\,.
\end{aligned}\end{align}}


\noindent
\textbf{2. Adiabatic integral.}
An alternative form is obtained through an adiabatic integral.
Introduce
{\small\begin{align}\begin{aligned}
    K_\lambda:=K_0+\lambda V\,,
    \quad 
    Q_\lambda:=\operatorname{Tr}\big\{e^{-\beta K_\lambda}\big\}\,,
    \quad
    \Omega_\lambda:=-\tfrac{1}{\beta}\log Q_\lambda\,,
    \quad 
    \lambda&\in[0,1]\,\,.
\end{aligned}\end{align}}
The Lie--Trotter formula and cyclicity of the trace give
{\small\begin{align}\begin{aligned}
    \tfrac{\partial\Omega_\lambda}{\partial\lambda}
    &=
    \tfrac{\partial}{\partial\lambda}
    \Big[
        -\tfrac{1}{\beta}
        \log\operatorname{Tr}\big\{e^{-\beta K_\lambda}\big\}
    \Big]
=
    \tfrac{1}{Q_\lambda}
    \operatorname{Tr}\big\{
        V e^{-\beta K_\lambda}
    \big\}
    =
    \langle V\rangle_{K_\lambda}\,.
\end{aligned}\end{align}}
Therefore, we find the \emph{adiabatic} representation of the \textbf{grand potential difference}:
{\small\begin{align}\begin{aligned}
    \Omega-\Omega_0
    &=
    \int_0^1
        \tfrac{\partial\Omega_\lambda}{\partial\lambda}\,d\lambda
    =
    \int_0^1
        \langle V\rangle_{K_\lambda}\,d\lambda\,.
\end{aligned}\end{align}}
This expression is used in superconductivity to compute the offset between the standard and superconducting phase.
\end{mytheorem}


%---------------------------------------
%================================================
\section{Functional formalism}
%================================================
%---------------------------------------------






\end{document}
